\section{A doubled seed and the common class}\label{sec:common}
The words of Section~\ref{sec:positive} need not avoid $110$. For example, seed $(0,1)$ and blocks $(1)$, $(0,2)$ produce
\[
 (0,1,1,0,2),
\]
a legal $000$- and $100$-avoider containing $110$. The following transform obtains a common-class lower bound with an explicit additional length cost.

\begin{lemma}\label{lem:doubling}
Replace the seed $(0,1,\ldots,r-1)$ of a word constructed in Section~\ref{sec:positive} by two consecutive copies of that seed, and increase every suffix label by $r$. The transformed word belongs to $\Av(000,100,110)$. For fixed $r$, this transform is injective and adds $r$ to the length.
\end{lemma}
\begin{proof}
The doubled seed is itself legal and has $2r-2$ ascents. The first original suffix entry is at most $r-1$, by the first-block order-statistic bound, so its entering edge from the original seed is a nonascent. Its transformed value lies between $r$ and $2r-1$, and is legal after the doubled seed; its entering edge is an ascent. Before that entry, the transformed ascent count exceeds the original count by $r-1$. After that entry, the difference is $r$. All later comparisons within the suffix are unchanged by translation, while every later suffix value has increased by $r$. Their legality therefore follows from the original legality bounds. This argument includes $r=1$.

The only repeated labels are the seed labels, each occurring exactly twice, so $000$ is absent. A seed label's first occurrence precedes every larger label, excluding $100$. After its second occurrence, all remaining seed entries are larger, and all suffix entries are at least $r$. No smaller entry can follow the two equal copies, excluding $110$. Finally, remove the first $r$ entries and subtract $r$ from the suffix to recover the original word; fixed $r$ makes the transform injective.
\end{proof}
The use of a doubled increasing seed has antecedents in Conway et al.\ \cite[Section~7]{conway} and in Report~241. Here its translated distinct suffix is combined with the triangular repair argument. We do not claim the doubled-seed idea itself as new.

Repeating all three fiber estimates from Proposition~\ref{prop:repair}, and padding with fresh records after the transform, gives the finite inequality
\begin{equation}\label{eq:commonfinite}
 c_{n+r+Q}\ge
 \frac{T(n,r)}{2^{r+1}(Q+1)\binom{n+Q}{r}}
 \quad(n>r\ge1),
\end{equation}
with exactly the same $Q$ as in \eqref{eq:Q}. The transformed lengths are $n+r+t$, so padding still has at most $Q+1$ possible original lengths. The rank and decoration fibers are unchanged by the injective transform.

To reach an arbitrary sufficiently large target $N$, set $L=\log N$ and choose
\begin{equation}\label{eq:commonparameters}
 r=\lfloor N/L\rfloor,\qquad
 n=N-r-\lceil4N/L^2\rceil,\qquad q=n-r.
\end{equation}
Now $M\sim N^2/(2L^2)$, $q=N-2N/L+O(N/L^2)$, and $Q\sim N/L^2$. Therefore $n+r+Q\le N$ eventually. The denominator in \eqref{eq:commonfinite} has logarithm $O(N\log L/L)$, by the same binomial estimate as before.

For clarity, the linear constant can be calculated directly, without applying a saddle formula outside its chosen dimension. We have
\[
 \log M=2L-2\log L-\log2+O(L/N),\qquad
 \log q=L-2/L+O(1/L^2).
\]
Use \eqref{eq:product} and the factorial estimate, with $q^2/M=O(L^2)$. Then
\begin{align*}
 \log T(n,r)
 &=q(\log M-\log q+1)+O(L^2+L)\\
 &=N(L-2\log L-\log2-1)+O(N\log L/L).
\end{align*}
The $-2N$ contribution from $q=N-2N/L+O(N/L^2)$ combines with the factorial's $+N$ contribution to give the $-1$ in $C_-$. Fresh-record padding and \eqref{eq:commonfinite} imply
\begin{equation}\label{eq:commonlower}
 \log c_N\ge N(\log N-2\log\log N+C_-)
                   -O(N\log\log N/\log N).
\end{equation}
Since $c_N\le d_N$, this lower bound also holds for $d_N$. Together with \eqref{eq:upperfinal}, it proves \eqref{eq:bracket} and completes Theorem~\ref{thm:main}.

The interval $[C_-,C_+]$ has width $2\log2$. The proof does not show that either $\kappa_d(N)$ or $\kappa_c(N)$ converges, nor that their possible limits agree. Improving the common-class lower construction or finding a matching $110$ comparison remains a separate problem.
